\begin{solution}{86}
    \begin{subsolution}
        \begin{subsubsolution}
            \begin{subsubsubsolution}
                \begin{subsubsubsubsolution}
                    对于能量为 $\frac{hc}{\lambda_0}$ 的光子，根据质能关系有
                    \begin{equation}
                        \frac{hc}{\lambda_{0}} = m c^{2}
                        \label{eq:1.s86}
                    \end{equation}
                    式中，$m$ 是该光子的等效质量。设在星球的引力场中，星球表面的重力加速度为 $g'$，根据万有引力定律有
                    \begin{equation}
                        g' = \frac{G M}{R^{2}}
                        \label{eq:2.s86}
                    \end{equation}
                    由能量守恒知，光子在 $A$、$B$ 两点的能量相等，即
                    \begin{equation}
                        \frac{hc}{\lambda_{0}} = \frac{hc}{\lambda'} + m g' L
                        \label{eq:3.s86}
                    \end{equation}
                    由\eqref{eq:1.s86}、\eqref{eq:3.s86}式得
                    \begin{equation}
                        \frac{1}{\lambda'} = \frac{1}{\lambda_{0}} \left(1 - \frac{g' L}{c^{2}}\right)
                        \label{eq:4.s86}
                    \end{equation}
                    由\eqref{eq:4.s86}式解得
                    \begin{equation}
                        \lambda' = \lambda_{0} \frac{1}{1 - \frac{g' L}{c^{2}}} \approx \lambda_{0} \left(1 + \frac{g' L}{c^{2}}\right)
                        \label{eq:5.s86}
                    \end{equation}
                    将\eqref{eq:2.s86}式代入\eqref{eq:5.s86}式得
                    \begin{equation}
                        \lambda' \approx \lambda_{0} \left(1 + \frac{G M}{R^{2}} \frac{L}{c^{2}}\right)
                        \label{eq:6.s86}
                    \end{equation}
                \end{subsubsubsubsolution}
                \begin{subsubsubsubsolution}
                    光在真空中的速度大小 $c$ 与发射器的运动状态无关。设 $A$ 发出的光以光速 $c$ 射向接收器 $B$ 直至到达 $B$ 的运动过程需要的时间为 $t$，有
                    \begin{equation}
                        c t = L + \frac{1}{2} a t^{2}
                        \label{eq:7.s86}
                    \end{equation}
                    设接收器 $B$ 接收到激光时相对于地面的速度为
                    \begin{equation}
                        v = a t
                        \label{eq:8.s86}
                    \end{equation}
                    由\eqref{eq:7.s86}、\eqref{eq:8.s86}式得
                    \begin{equation}
                        v \approx \frac{a L}{c}
                        \label{eq:9.s86}
                    \end{equation}
                    已略去 $\frac{aL}{c^{2}}$ 的高阶小量。按照多普勒频移公式，$B$ 点接收到的光的波长 $\lambda''$ 满足
                    \begin{equation}
                        \frac{c}{\lambda''} = \sqrt{\frac{1-\beta}{1+\beta}} \frac{c}{\lambda_{0}}
                        \label{eq:10.s86}
                    \end{equation}
                    式中
                    \begin{equation}
                        \beta = \frac{v}{c} = \frac{a L}{c^{2}}
                        \label{eq:11.s86}
                    \end{equation}
                    由于 $\beta \ll 1$，有
                    \begin{equation}
                        \sqrt{\frac{1-\beta}{1+\beta}} = \sqrt{1 - \frac{2\beta}{1+\beta}} \approx 1 - \frac{\beta}{1+\beta} = \frac{1}{1+\beta}
                        \label{eq:12.s86}
                    \end{equation}
                    将\eqref{eq:11.s86}、\eqref{eq:12.s86}式代入\eqref{eq:10.s86}式得，$B$ 点的接收器收到的激光波长为
                    \begin{equation}
                        \lambda'' = \lambda_{0} \left(1 + a \frac{L}{c^{2}}\right)
                        \label{eq:13.s86}
                    \end{equation}
                \end{subsubsubsubsolution}
                \begin{subsubsubsubsolution}
                    根据等效性原理，若将在（i）和（ii）中所计算的 $B$ 点接收器收到的激光的波长 $\lambda'$ 和 $\lambda''$ 视为同一事件，\eqref{eq:6.s86}式和\eqref{eq:13.s86}式便应表达相同的物理规律；比较（i）和（ii）的结果有
                    \begin{equation}
                        a = \frac{G M}{R^{2}}
                        \label{eq:14.s86}
                    \end{equation}
                \end{subsubsubsubsolution}
            \end{subsubsubsolution}
        \end{subsubsolution}
        \begin{subsubsolution}
            在地球表面，重力加速度大小为 $g$，用 $g$ 代替\eqref{eq:4.s86}式中的 $g'$ 得，$B$ 点接收到的光子的频率 $\nu'$ 为
            \begin{equation}
                \nu' = \nu_{0} \left(1 - \frac{g L}{c^{2}}\right)
                \label{eq:15.s86}
            \end{equation}
            由\eqref{eq:15.s86}式得，地球重力作用造成的光子频移为
            \[
                \Delta \nu = \nu' - \nu_{0} = - \nu_{0} \frac{g L}{c^{2}}
            \]
            可见，重力作用造成的频率减小了 $\nu_{0}\frac{gL}{c^{2}}$（红移），需要由穆斯堡尔探测器的相对于地面的运动来进行补偿，以便用穆斯堡尔探测器能探测到该光子。设穆斯堡尔探测器相对于地面沿 $BA$ 方向运动的速度为 $v$ 时，穆斯堡尔探测器能探测到的光子的频率。按照多普勒频移公式有
            \begin{equation}
                \nu_{0} = \sqrt{\frac{1+\beta}{1-\beta}} \nu'
                \label{eq:16.s86}
            \end{equation}
            由\eqref{eq:12.s86}、\eqref{eq:16.s86}式可知
            \begin{equation}
                \nu' = \frac{\nu_{0}}{1+\beta} \approx \nu_{0} \left(1 - \frac{v}{c}\right)
                \label{eq:17.s86}
            \end{equation}
            由\eqref{eq:15.s86}、\eqref{eq:17.s86}式得
            \begin{equation}
                \frac{g L}{c^{2}} = \frac{v}{c}
                \label{eq:18.s86}
            \end{equation}
            即
            \begin{equation}
                v = \frac{g L}{c}
                \label{eq:19.s86}
            \end{equation}
            由\eqref{eq:19.s86}式和题给数据得
            \begin{equation}
                v = \frac{9.81 \times 22.6}{3.00 \times 10^{8}}\,\mathrm{m\cdot s^{-1}} = 7.39 \times 10^{-7}\,\mathrm{m\cdot s^{-1}}
                \label{eq:20.s86}
            \end{equation}
        \end{subsubsolution}
    \end{subsolution}
\end{solution}